%! Author = sriadityavedantam
%! Date = 4/8/26

\section{Dimension of Divisors}

Let $X$ be a smooth projective variety and let $D$ be a divisor on $X$.
Define
\[
    \Ell(D) = \{0\} \cup \{f \in \c(X) : \div(f) + D \geq 0\}.
\]
Then $\Ell(D)$ is a $\c$-vector space of finite dimension.

\begin{definition}[Dimension of a Divisor]
    The dimension of a divisor $D$ is
    \[
        \ell(D) \coloneqq \dim_\c \Ell(D).
    \]
\end{definition}

\subsection*{Basic Bound}

\begin{theorem*}{
    Let $X$ be a smooth projective curve and $D$ a divisor on $X$.
    If $D$ is effective or a $0$-divisor, then
    \[
        \ell(D) \leq \deg(D) + 1.
    \]
}
\end{theorem*}

If $D = \sum a_i p_i$, then $\deg(D) = \sum a_i$.
In particular, this theorem implies $\ell(D)$ is finite.
If $D = 0$, then
\[
    \Ell(D) = \c[X] = \c,
\]
so $\ell(D) = 1$.
Since $\deg(D) = 0$, we obtain equality:
\[
    \ell(D) = \deg(D) + 1.
\]

\subsection*{Case $X = \pp^1$}

If $D \sim D'$, then $\Ell(D) \cong \Ell(D')$.
Thus we may assume
\[
    D = \deg(D)\cdot p_\alpha.
\]

Then
\[
    \Ell(D)
    = \{0\} \cup \{f \in \c(\pp^1) : \div(f) + D \geq 0\}.
\]

This means $f$ is regular on $\pp^1 \setminus \{p_\alpha\}$ and has a pole of order at most $\deg(D)$ at $p_\alpha$.
Choose an affine chart
\[
    \pp^1 \setminus \{p_\alpha\} \cong \aa^1_t.
\]
Then $f \in \c[t]$.
Using the coordinate $u = \frac{1}{t}$ near $p_\alpha$, the condition becomes
\[
    \deg(f) \leq \deg(D).
\]

Hence
\[
    \Ell(D)
    \cong \{f \in \c[t] : \deg(f) \leq \deg(D)\}
    \cong \c^{\deg(D)+1}.
\]

Therefore
\[
    \ell(D) = \deg(D) + 1.
\]

\subsection*{Nontrivial Case: $X \neq \pp^1$}

Let $X$ be a smooth projective curve not isomorphic to $\pp^1$, and let
\[
    D = p \in X.
\]

Then
\[
    \Ell(D)
    = \{0\} \cup \{f \in \c(X) : \div(f) + p \geq 0\}.
\]

Thus $f$ is regular on $X \setminus \{p\}$ and has at most a first-order pole at $p$.

\begin{proof}
    If $f$ has no poles, then $f \in \c[X]$, hence $f$ is constant.

    Assume $f$ has a first-order pole at $p$.
    Then $f$ defines a morphism
    \[
        f : X \to \pp^1.
    \]

    Since the pole has order $1$, we have $\deg(f) = 1$.
    Thus $f$ is an isomorphism, implying $X \cong \pp^1$, a contradiction.

    Therefore $f$ must be constant.
    Hence
    \[
        \Ell(D) \cong \c,
        \quad
        \ell(D) = 1.
    \]

    Since $\deg(D) = 1$, we obtain
    \[
        \ell(D) = 1 < \deg(D) + 1 = 2.
    \]
\end{proof}

This shows that equality $\ell(D) = \deg(D)+1$ does not hold in general.